
\documentclass{../concepts/prepareprotokol} %všechny balíčky pro protokol + titulní strana
\usepackage{graphicx}
\usepackage{newfloat}
\usepackage{caption}
\usepackage{booktabs}
\usepackage{amsmath}
\usepackage{newfloat}
\usepackage{caption}
\usepackage{siunitx}
\sisetup{inter-unit-product =$\cdot$}
\sisetup{separate-uncertainty=true}

\praktikum{II} %číslo praktika (I, II nebo III)
\autor{David Němec} %jméno autora
\datum{8.\,12.\,2025} %datum měření
\cislo{7} %číslo úlohy
\nazev{Měření indukčnosti a kapacity metodou přímou} %název úlohy

\DeclareFloatingEnvironment[fileext=los, placement={htbp}, name=Obrázek]{schema}

\newcommand{\nobars}{
Chybové úsečky nebyly pro svou malou velikost pro přehlednost vykresleny.}


\begin{document}
\renewcommand{\figurename}{Graf}
\renewcommand{\refname}{Literatura}

\maketitle

% ========================== PRACOVNÍ ÚKOL ============================
\section{Pracovní úkol}
\begin{enumerate}
    \item Změřte velikost kapacit kondenzátorů z kapacitní dekády.
    \item Změřte závislost indukčnosti cívky na procházejícím proudu pro tyto případy:
    \begin{enumerate}
        \item cívka bez jádra (0-250 mA)
        \item cívka s otevřeným jádrem (0-250 mA)
        \item cívka s uzavřeným jádrem (0-250 mA)
    \end{enumerate}

    \item Změřte odpor cívky multimetrem a určete její kvalitu.
    \item Závislosti indukčnosti cívky na proudu měřené v úkolu 2 zakreslete do grafu přímo v praktiku.
\end{enumerate}

% ============================= TEORIE ================================
\section{Teorie}
\subsection{Odpor a kapacita kondenzátoru}

Reálný kondenzátor, který má kromě své kapacity $C$ i nenulový vlastní odpor $R_C$, lze popsat náhradním schématem (obr. \ref{sch:kondenzator}), kdy jsou paralelně zapojeny ideální kondenzátor $C$ a rezistor $R_C$. Kondenzátorem teče posuvný proud $i_p$ a rezistorem vodivostní proud $i_v$.

Pokud zapojíme reálný kondenzátor ke zdroji střídavého napětí o úhlové frekvenci $\omega$ a pro jeho odpor bude platit $R_C \gg (\omega C)^{-1}$, lze kapacitu kondenzátoru určit ze vztahu \cite{studijni}
\begin{equation}
    C = \frac{1}{2 \pi f} \frac{I}{U} \, ,
\end{equation}
kde $f$ je frekvence střídavého napětí, $I$ a $U$ jsou efektivní hodnoty proudu resp. napětí.

Při připojení ke stejnosměrnému napětí se kapacita kondenzátoru neprojeví a pro jeho odpor bude platit Ohmův zákon \cite{studijni4}
\begin{equation}
    R_C = \frac{U}{I} \, .
\end{equation}

\subsection{Odpor a indukčnost cívky}
Reálnou cívku s indukčností $L$ a odporem $R_L$ lze popsat náhradním schématem (obr. \ref{sch:civka}), kdy jsou sériově zapojeny ideální cívka $L$ a rezistor $R_L$.

Pokud zapojíme cívku ke zdroji střídavého napětí o frekvenci $f$, lze indukčnost cívky vypočíst ze vztahu \cite{studijni}
\begin{equation}
    L = \frac{\sqrt{\frac{U^2}{I^2}-R_L^2}}{2 \pi f} \, .
\end{equation}

Měřítkem kvality cívky je její činitel jakosti, definovaný jako (upraveno podle \cite{studijni})
\begin{equation}
    Q = \frac{2 \pi f L}{R_L} \, .
\end{equation}

\subsection{Čtyřbodové zapojení rezistoru}

Pro přesnější měření odporu rezistoru (ale i cívky) lze využít čtyřbodového zapojení podle schématu na obrázku \ref{sch:four}. Pokud je voltmetr připojen k napěťovým svorkám $b$ a $b'$, nebude měření ovlivněno spádem napětí na přívodních vodičích mezi rezistorem a proudovými svorkami $a$ a $a'$. Za předpokladu nízkých hodnot $R$ oproti odporu voltmetru poteče přívodními vodiči k napěťovým svorkám zanedbatelný proudu a ani odpor těchto vodičů se neuplatní \cite{studijni4}.

\subsection{Metoda měření}

\begin{schema}[htbp]
    \centering
    \includegraphics[width=0.5\linewidth]{uloha4/fourpoints.png}
    \caption{Schema yap}
    \label{sch:4}
\end{schema}

\begin{figure}[htbp]
    \centering
    \includegraphics[width=0.5\linewidth]{uloha27/schema.png}
    \caption{Schéma zapojení [LITERATURA]}
    \label{fig:schema}
\end{figure}



\subsection{Měřicí přístroje a jejich nejistoty}
\begin{enumerate}
    \item PŘÍSTROJ:
    
    POZNÁMKY

\end{enumerate}


% ======================== VÝSLEDKY MĚŘENÍ ===========================
\section{Výsledky měření}


\begin{figure}
    \centering
    \includegraphics[width=0.5\linewidth]{uloha27/GRAF1.png}
    \caption{Závislost kapacity deskového kondenzátoru na napětí.}
    \label{fig:osciloskop}
\end{figure}


% ============================ DISKUZE ===============================
\section{Diskuze}


% ============================= ZÁVĚR ================================
\section{Závěr}


% =========================== LITERATURA =============================
\begin{thebibliography}{9}

\bibitem{studijni} Kolektiv ZFP KVOF MFF UK: Měření indukčnosti a kapacity metodou přímou - studijní text [online]. [cit. 12.12.2025], dostupné z https://physics.mff.cuni.cz/vyuka/zfp/\_media/zadani/texty/txt\_207.pdf

\bibitem{studijni4} Kolektiv ZFP KVOF MFF UK: Měření malých odporů [online]. [cit. 12.12.2025], dostupné z~https://physics.mff.cuni.cz/vyuka/zfp/\_media/zadani/texty/txt\_204.pdf

\bibitem{pokyny} Kolektiv ZFP KVOF MFF UK: Měření indukčnosti a kapacity metodou přímou - pokyny k měření [online]. [cit. 12.12.2025], dostupné \\
z https://physics.mff.cuni.cz/vyuka/zfp/\_media/zadani/pokyny/mereni\_207.pdf

\bibitem{englich} J. Englich: Úvod do praktické fyziky I: Zpracování výsledků měření. 1. vyd. Praha: Matfyzpress, 2006

\bibitem{mikulcak} Mikulčák a kolektiv: Matematické, fyzikální a chemické tabulky pro střední školy: relativní permitivita, Praha: Státní pedagogické nakladatelství, n. p., 1988

\end{thebibliography}
\end{document}